\section{Hopf structure, primitive elements, and why the logarithm is a Lie series}\label{sec:hopf}
We prove the primitive-element characterization used by the source page
rather than treating it as a black box. The proof simultaneously constructs
the projection that produces Dynkin's formula.

\subsection{The coproduct must take values in a completion}
The relevant target of the coproduct on infinite series is the
\emph{completed} tensor product
\[
 \widehat T\ot\widehat T
 =\prod_{n\geq0}\bigoplus_{p+q=n}V^{\otimes p}\otimes V^{\otimes q},
\]
with multiplication and infinite sums interpreted by total degree exactly as
in $\widehat T$. The ordinary algebraic tensor product
$\widehat T\otimes\widehat T$ is generally too small: we show in
\cref{app:pbw} that already the coproduct of the one-variable series
$\sum_{n\geq0}X^n$ is not a finite sum of separated tensors. Define the
continuous algebra homomorphism
\begin{equation}\label{eq:coproduct}
 \Delta:\widehat T\longrightarrow\widehat T\ot\widehat T,
 \qquad
 \Delta X=X\otimes1+1\otimes X,\quad
 \Delta Y=Y\otimes1+1\otimes Y.
\end{equation}
Since $\widehat T$ is the completion of the free algebra on $X,Y$, an
algebra homomorphism on $T(V)$ is determined freely by the images of $X$ and
$Y$; and because $\Delta$ neither raises nor lowers total degree, it
extends uniquely by continuity to $\widehat T$. For a word
$w=a_1\cdots a_n$ and a subset $I\subseteq\{1,\ldots,n\}$ let $w_I$ be the
subword formed by the letters at the positions in $I$, in their original
order, with $w_\varnothing=1$. Then
\begin{equation}\label{eq:unshuffle}
 \Delta w=\sum_{I\subseteq\{1,\ldots,n\}}w_I\otimes w_{I^c}.
\end{equation}
This \emph{unshuffle} formula follows by multiplying out the $n$ factors
$a_j\otimes1+1\otimes a_j$: each choice of one term from each factor
assigns every letter to the left or right tensor factor, the letters
assigned to the left concatenate in their original order to $w_I$, and those
assigned to the right concatenate to $w_{I^c}$. Coassociativity,
$(\Delta\otimes\id)\Delta=(\id\otimes\Delta)\Delta$, follows in the same way:
either iterated coproduct assigns each letter independently to one of three
tensor factors, and the resulting sum over assignments does not depend on
the order in which the two splittings are performed. The counit
$\eps:\widehat T\to\kk$ extracts the constant term; the identities
$(\eps\otimes\id)\Delta=\id=(\id\otimes\eps)\Delta$ hold on words by
\eqref{eq:unshuffle}, since the only term with $w_I=1$ is $I=\varnothing$.

Define the antipode on words by
\begin{equation}\label{eq:antipode}
 S(a_1\cdots a_n)=(-1)^n a_n\cdots a_1,\qquad S(1)=1,
\end{equation}
extended linearly and continuously. It is an antiautomorphism:
$S(uv)=S(v)S(u)$, since reversing a concatenation reverses the factors in
the opposite order and the signs multiply. For linear maps $f,g$ on
$\widehat T$ the \emph{convolution} is $f*g=m(f\otimes g)\Delta$, where $m$
is multiplication; by coassociativity convolution is associative, with unit
$\eta\eps$, where $\eta$ inserts scalars.

\begin{lemma}[Antipode identities]\label{lem:antipode}
$\id*S=S*\id=\eta\eps$, that is, for every word $w$,
\begin{equation}\label{eq:antipodeidentities}
 m(\id\otimes S)\Delta w=m(S\otimes\id)\Delta w=\eps(w)1.
\end{equation}
\end{lemma}
\begin{proof}
Write $\Delta v=\sum v_{(1)}\otimes v_{(2)}$ in Sweedler notation for the
finite sum \eqref{eq:unshuffle} and put $K(v)=\sum v_{(1)}S(v_{(2)})$. For a letter $a$
and a word $v$, the unshuffle formula for $av$ splits according to whether
the first letter is assigned to the left or the right factor:
\[
 \Delta(av)=\sum av_{(1)}\otimes v_{(2)}+\sum v_{(1)}\otimes av_{(2)}.
\]
Using $S(av_{(2)})=-S(v_{(2)})a$ we obtain
\[
 K(av)=a\sum v_{(1)}S(v_{(2)})-\sum v_{(1)}S(v_{(2)})a=aK(v)-K(v)a.
\]
Since $K(1)=1\cdot S(1)=1$, we get $K(a)=a-a=0$ for every letter, and then
inductively $K(w)=0$ for every nonempty word. This proves the first identity.
For the second, put $K'(v)=\sum S(v_{(1)})v_{(2)}$ and split the coproduct of
$va$ according to the position of the last letter $a$:
$\Delta(va)=\sum v_{(1)}a\otimes v_{(2)}+\sum v_{(1)}\otimes v_{(2)}a$, whence
$K'(va)=-aK'(v)+K'(v)a$, and the same induction applies.
\end{proof}
Thus $T(V)$, and by continuity its completion, carries the asserted Hopf
algebra structure, with no appeal to an external structure theorem.

\begin{definition}
A series $P\in\widehat T$ is \emph{primitive} if
$\Delta P=P\otimes1+1\otimes P$. A series $G\in1+F^1\widehat T$ is
\emph{group-like} if $\Delta G=G\otimes G$.
\end{definition}
The normalization of the constant term is part of the definition of a
group-like element: the zero series satisfies $\Delta0=0\otimes0$ but does not
belong to a group. A primitive element has zero constant term, since a
primitive scalar $c$ would satisfy $c\otimes1=2c\otimes1$.

\begin{lemma}[Exponentials of primitives]\label{lem:grouplike}
Let $P\in F^1\widehat T$. Then $P$ is primitive if and only if $e^P$ is
group-like. The group-like elements form a group under multiplication.
\end{lemma}
\begin{proof}
Since $\Delta$ is a continuous algebra homomorphism, it commutes with the
exponential series: $\Delta(e^P)=e^{\Delta P}$. If $P$ is primitive, then
$\Delta P=P\otimes1+1\otimes P$ is a sum of two commuting elements, because
$(P\otimes1)(1\otimes P)=P\otimes P=(1\otimes P)(P\otimes1)$. By
\cref{lem:explog} applied in $\widehat T\ot\widehat T$,
\[
 \Delta(e^P)=e^{P\otimes1+1\otimes P}=e^{P\otimes1}e^{1\otimes P}
 =(e^P\otimes1)(1\otimes e^P)=e^P\otimes e^P.
\]
Conversely, if $G=e^P$ is group-like, then
$\Delta G=G\otimes G=(G\otimes1)(1\otimes G)$ is a product of two commuting
elements of $1+F^1(\widehat T\ot\widehat T)$, and the commuting-logarithm
rule of \cref{lem:explog} gives
\[
 \Delta P=\log(\Delta G)=\log(G\otimes1)+\log(1\otimes G)
 =P\otimes1+1\otimes P,
\]
where $\log(G\otimes1)=(\log G)\otimes1$ because $A\mapsto A\otimes1$ is a
continuous algebra homomorphism. Finally, if $G,H$ are group-like then
$\Delta(GH)=(G\otimes G)(H\otimes H)=GH\otimes GH$, and $G$ is invertible by
its geometric series, with
$\Delta(G^{-1})=(\Delta G)^{-1}=(G\otimes G)^{-1}=G^{-1}\otimes G^{-1}$.
\end{proof}

\subsection{An explicit projection: the Dynkin--Specht--Wever identity}
Let $N$ be the degree operator, $N(w)=|w|\,w$ and $N(1)=0$. It is a
derivation of $T(V)$, because degrees add under concatenation. Consider the
convolution
\begin{equation}\label{eq:dynkinoperator}
 D=N*S=m(N\otimes S)\Delta.
\end{equation}
The order of the two convolution factors is what produces
\emph{right}-nested rather than left-nested brackets. This convolution
operator is a special case of the generalized Dynkin operators studied in
\cite{menouspatras}; the elementary computation below contains the entire
argument required here.

\begin{lemma}[Right Dynkin operator]\label{lem:dsw}
For every nonempty word $w$, $D(w)=R(w)$. If $P$ is a homogeneous primitive
element of degree $n\geq1$, then
\begin{equation}\label{eq:dsw}
 R(P)=D(P)=nP,\qquad\text{that is,}\qquad
 P=\frac1n\sum_{|w|=n}\coeff{w}P\,R(w).
\end{equation}
\end{lemma}
\begin{proof}
From the definitions, $D(1)=N(1)S(1)=0$ and, for a letter $a$,
$D(a)=N(a)S(1)+N(1)S(a)=a$. Let $a$ be a letter and $w$ a nonempty word, and
split $\Delta(aw)$ as in the proof of \cref{lem:antipode}:
$\Delta(aw)=\sum aw_{(1)}\otimes w_{(2)}+\sum w_{(1)}\otimes aw_{(2)}$.
Using $N(aw_{(1)})=aw_{(1)}+aN(w_{(1)})$ (derivation property with
$N(a)=a$) and $S(aw_{(2)})=-S(w_{(2)})a$, we get
\begin{align*}
 D(aw)
 &=\sum aw_{(1)}S(w_{(2)})+\sum aN(w_{(1)})S(w_{(2)})
   -\sum N(w_{(1)})S(w_{(2)})a\\
 &=aK(w)+aD(w)-D(w)a,
\end{align*}
where $K(w)=\sum w_{(1)}S(w_{(2)})=\eps(w)1=0$ for nonempty $w$ by
\cref{lem:antipode}. Hence $D(aw)=[a,D(w)]$, and induction on the length
starting from $D(a)=a$ gives $D(a_1\cdots a_n)=[a_1,D(a_2\cdots a_n)]
=\cdots=R(a_1\cdots a_n)$.

If $P$ is homogeneous primitive of degree $n$, then its coproduct has only
the two terms $P\otimes1$ and $1\otimes P$, and
\[
 D(P)=m(N\otimes S)(P\otimes1+1\otimes P)=N(P)S(1)+N(1)S(P)=nP+0.
\]
Since $D$ is linear and $D(w)=R(w)$ on words, $D(P)=\sum_w\coeff{w}P\,R(w)$,
which gives the second form of \eqref{eq:dsw} after division by $n$.
\end{proof}

Let $\cL=\cL(X,Y)\subseteq T(V)$ be the Lie subalgebra generated by $X,Y$
under the commutator bracket, and $\cL_n=\cL\cap V^{\otimes n}$. Since
$X,Y$ are homogeneous, $\cL$ is spanned by homogeneous bracket monomials, so
$\cL=\bigoplus_n\cL_n$; we write $\widehat\cL=\prod_{n\geq1}\cL_n$.

\begin{theorem}[Friedrichs' characterization; Dynkin--Specht--Wever]\label{thm:friedrichs}
A homogeneous element $P\in V^{\otimes n}$, $n\geq1$, is primitive if and
only if $P\in\cL_n$. For such $P$, \eqref{eq:dsw} holds, so $D/n$ is an
idempotent projection of $V^{\otimes n}$ onto $\cL_n$ whose kernel consists
of the degree-$n$ elements annihilated by $D$. The primitive elements of
$\widehat T$ are exactly the elements of $\widehat\cL$, the formal
degreewise sums of Lie polynomials:
\begin{equation}\label{eq:friedrichs}
 \Prim(\widehat T)=\prod_{n\geq1}\cL_n.
\end{equation}
\end{theorem}
\begin{proof}
The generators $X,Y$ are primitive by definition of $\Delta$. If $P,Q$ are
primitive, then
\[
 \Delta[P,Q]=[\Delta P,\Delta Q]
 =[P\otimes1+1\otimes P,\ Q\otimes1+1\otimes Q]
 =[P,Q]\otimes1+1\otimes[P,Q],
\]
because the mixed terms cancel: $(P\otimes1)(1\otimes Q)=P\otimes Q
=(1\otimes Q)(P\otimes1)$, and similarly with $P$ and $Q$ interchanged.
Hence every bracket monomial in $X,Y$ is primitive, every Lie polynomial is
primitive by linearity, and every degreewise-convergent sum of Lie
polynomials is primitive by continuity of $\Delta$.

Conversely, let $P\in\widehat T$ be primitive. Its constant term vanishes,
as noted above. Since $\Delta$ preserves total degree, comparing the
degree-$n$ parts of $\Delta P=P\otimes1+1\otimes P$ shows that each
homogeneous component $P_n$ is primitive. By \cref{lem:dsw},
$P_n=\frac1n\sum_{|w|=n}\coeff{w}P_n\,R(w)$ is a linear combination of
right-nested brackets, hence an element of $\cL_n$. Division by $n$ is
exactly where characteristic zero is used. This proves the homogeneous
characterization and \eqref{eq:friedrichs}.

Finally, since $D$ always takes values in $\cL$ (its value on any word is a
nested bracket) and every element of $\cL_n$ is primitive with $D(P)=nP$,
we have $D(D(Q))=nD(Q)$ for every $Q\in V^{\otimes n}$; so $D/n$ is
idempotent with image $\cL_n$.
\end{proof}

\begin{theorem}[Formal BCH theorem]\label{thm:formalbch}
There is a unique $Z\in F^1\widehat T$ with $e^Z=e^Xe^Y$. Every homogeneous
component $Z_n$ is a Lie polynomial with rational coefficients, $Z_1=X+Y$,
and $Z-X-Y$ belongs degreewise to the Lie ideal of $\cL$ generated by
$[X,Y]$.
\end{theorem}
\begin{proof}
Existence and uniqueness follow from \cref{lem:explog}. Since $X$ and $Y$
are primitive, $e^X$ and $e^Y$ are group-like by \cref{lem:grouplike}, their
product $e^Xe^Y$ is group-like because group-like elements form a group, and
$Z=\log(e^Xe^Y)$ is primitive by \cref{lem:grouplike} again. By
\cref{thm:friedrichs} each $Z_n$ lies in $\cL_n$. Rationality follows from
\eqref{eq:wordcut}, whose ingredients are factorials and the coefficients
$(-1)^{k-1}/k$. Taking degree one in \eqref{eq:blocks} gives $Z_1=X+Y$.

For the ideal assertion, let $J$ be the Lie ideal of $\cL$ generated by
$[X,Y]$. A right-nested bracket $R(w)$ with $|w|\geq2$ has innermost
bracket $[a_{n-1},a_n]$, which is either zero or $\pm[X,Y]$, and an iterated
bracket one of whose factors lies in the ideal $J$ lies in $J$; since each
$Z_n$ with $n\geq2$ is a combination of such brackets by \eqref{eq:dsw}, it
lies in $J$. (Equivalently, in the quotient Lie algebra $\cL/J$ the two
generators commute, so $\cL/J$ is abelian and the image of $Z$ under the
Lie homomorphism $\cL\to\cL/J$ is $\bar X+\bar Y$ by \cref{prop:symmetry}
applied in an abelian Lie algebra, which again says $Z_n\in J$ for
$n\geq2$.)
\end{proof}
The phrase ``every term beyond $X+Y$ contains $[X,Y]$'' is thus an
ideal-theoretic statement; it does not mean that an arbitrary display of
$Z_n$ must visibly show $[X,Y]$ as the innermost bracket of every summand.

\subsection{Universality and enveloping algebras}\label{sec:universal}
A homogeneous Lie polynomial can be evaluated at any pair of elements of any
Lie algebra $\g$ over $\kk$. The identification of $\cL(X,Y)$ with the
abstract free Lie algebra, and the fact that the free associative algebra is
its universal enveloping algebra, rest on the injectivity of $\g\to U(\g)$,
which is proved in \cref{app:pbw} by the ordered Poincar\'e--Birkhoff--Witt
theorem (\cref{thm:pbw}). The argument is as follows. Given $x,y\in\g$, the
associative homomorphism $T(V)\to U(\g)$ determined by $X\mapsto x$,
$Y\mapsto y$ sends every iterated commutator of $X,Y$ to the corresponding
iterated bracket of $x,y$, which lies in the image of $\g$. Its restriction
to $\cL$ is therefore a Lie homomorphism $\cL\to\g$, and it is the unique
one sending $X,Y$ to $x,y$ because $X,Y$ generate $\cL$. This is exactly the
universal property of the free Lie algebra. Consequently an associative
polynomial identity between Lie polynomials in $X,Y$, verified as an
equality of word coefficients in $T(V)$, is an identity valid under
substitution in \emph{every} Lie algebra.

\begin{proposition}[Enveloping algebra of the free Lie algebra]\label{prop:enveloping}
The universal enveloping algebra of $\cL(X,Y)$ is canonically $T(V)$, and
the Hopf structure of \cref{sec:hopf} is the enveloping Hopf structure
determined by $\Delta u=u\otimes1+1\otimes u$, $\eps(u)=0$, $S(u)=-u$ for
$u\in\cL$. The same formulas define a Hopf algebra structure on $U(\g)$ for
every Lie algebra $\g$.
\end{proposition}
\begin{proof}
The inclusion $\cL\to T(V)$ is a Lie homomorphism into the commutator
algebra of $T(V)$, so by the universal property of $U(\cL)$ it extends to an
associative homomorphism $\alpha:U(\cL)\to T(V)$. The assignment
$X\mapsto X$, $Y\mapsto Y$ into $U(\cL)$ extends, by freeness of $T(V)$, to
an associative homomorphism $\beta:T(V)\to U(\cL)$. Both composites
$\alpha\beta$ and $\beta\alpha$ fix $X$ and $Y$; since $X,Y$ generate
$T(V)$ as an algebra and generate $U(\cL)$ as an algebra (because they
generate $\cL$ as a Lie algebra and $\cL$ generates $U(\cL)$), both
composites are identities. The Hopf assertion for general $\g$ is proved in
\cref{app:pbw}; on generators the two structures agree, and each is
determined by its values on generators.
\end{proof}
For an arbitrary Lie algebra $\g$ without a topology, the unambiguous
universal statement is
\[
 \BCH(tx,ty)=\sum_{n\geq1}t^nZ_n(x,y)\in t\,\g[[t]],
\]
obtained by evaluating in $U(\g)[[t]]$. An infinite sum at $t=1$ in the
untopologized vector space $\g$ has no meaning until nilpotence
(\cref{sec:nilpotent}), a complete filtration, or a convergence argument
(\cref{sec:convergence}) has been supplied.
